\begin{summarybox}
	\textbf{\textcolor{CSummary}{一句话核心}}

	\textbf{两圆交于 $A,B$ 时，公共弦所在直线垂直于连心线；先由方向定斜率，再用两个勾股方程相减求弦中点，最后写直线方程。}

	\medskip
	\textbf{\textcolor{CSummary}{使用条件}}
	\begin{itemize}[leftmargin=*, itemsep=0.5em, topsep=0.4em]
		\item \textbf{条件1（两点相交）：} 两圆圆心距 $d>0$，且 $|r_1-r_2|<d<r_1+r_2$。
		\item \textbf{条件2（圆心与半径已知）：} $O_1(a_1,b_1)$、$O_2(a_2,b_2)$，半径 $r_1,r_2$。
	\end{itemize}

	\medskip
	\textbf{\textcolor{CSummary}{快懂：四步几何法}}
	\begin{itemize}[leftmargin=*, itemsep=0.4em, topsep=0.3em]
		\item \textbf{定方向：} $AB\perp O_1O_2$。
		\item \textbf{定位置：} 两个勾股式相减，$\displaystyle t=\frac{d^2+r_1^2-r_2^2}{2d}$。
		\item \textbf{求中点：} $\displaystyle M=O_1+\frac{t}{d}(O_2-O_1)$。
		\item \textbf{写直线：} $\displaystyle (a_2-a_1)(x-x_M)+(b_2-b_1)(y-y_M)=0$。
	\end{itemize}

	\medskip
	\textbf{\textcolor{CSummary}{快记：代数一步法}}
	\[
		\boxed{\underbrace{(x^2+y^2+D_1x+E_1y+F_1)}_{C_1\text{ 左端}}-\underbrace{(x^2+y^2+D_2x+E_2y+F_2)}_{C_2\text{ 左端}}=0}
	\]
	即 $(D_1-D_2)x+(E_1-E_2)y+F_1-F_2=0$。

	\medskip
	\textbf{\textcolor{CSummary}{关键提醒}}
	\begin{itemize}[leftmargin=*, itemsep=0.5em, topsep=0.4em]
		\item \textbf{中点不等于两圆心中点：} 半径不同一般有 $t\ne d/2$；$M$ 甚至可以在圆心线段之外。
		\item \textbf{特殊方向不用斜率：} 连心线水平或竖直，直接写 $x=x_M$ 或 $y=y_M$。
		\item \textbf{不要倒推错误：} 相减所得是公共弦\textbf{所在直线}，不是两圆的全部公共点。
	\end{itemize}
\end{summarybox}
